\begin{example}[Tensor products and quotients]\label{ex:vi24-a01}
For an $A$-algebra $B$ and ideal $I\subseteq B$, scalar extension sends $B/I$ to $(B\otimes_AA')/I(B\otimes_AA')$.
\end{example}

\begin{example}[Tensor products and localization]\label{ex:vi24-a02}
The canonical map $B_b\otimes_AA'\to(B\otimes_AA')_{b\otimes1}$ is an isomorphism.
\end{example}

\begin{example}[Iterated scalar extension]\label{ex:vi24-a03}
For $A\to A'\to A''$, associativity gives $(B\otimes_AA')\otimes_{A'}A''\cong B\otimes_AA''$.
\end{example}

\begin{example}[Products after scalar extension]\label{ex:vi24-a04}
Associativity and symmetry of tensor products yield $(B\otimes_AC)\otimes_AA'\cong(B\otimes_AA')\otimes_{A'}(C\otimes_AA')$.
\end{example}

\begin{example}[Base change can disconnect]\label{ex:vi24-a05}
$\Spec\mathbb C$ is connected as an $\mathbb R$-scheme, but after base change to $\mathbb C$ it becomes $\Spec(\mathbb C\times\mathbb C)$, which is disconnected.
\end{example}

\begin{example}[Base change can change irreducibility]\label{ex:vi24-a06}
The real scheme $\Spec\mathbb R[x]/(x^2+1)$ is irreducible, but its complexification has two irreducible components.
\end{example}

\begin{example}[Purely inseparable base change can create nilpotents]\label{ex:vi24-a07}
In characteristic $p$, let $k=\mathbb F_p(s^p)$ and $L=\mathbb F_p(s)$. Then $L\otimes_kL\cong L[u]/((u-s)^p)$, which is nonreduced.
\end{example}

\begin{example}[Separable splitting]\label{ex:vi24-a08}
For a finite separable extension $L/k$ and an algebraic closure $\overline k$, $L\otimes_k\overline k$ is a product of copies of $\overline k$, one for each $k$-embedding $L\hookrightarrow\overline k$.
\end{example}

\begin{example}[Closed immersion stability]\label{ex:vi24-a09}
A surjection $B\twoheadrightarrow C$ remains surjective after tensoring with $A'$; therefore affine closed immersions survive arbitrary base change.
\end{example}

\begin{example}[Open immersion stability]\label{ex:vi24-a10}
Localization commutes with scalar extension, so principal open immersions survive base change; affine localization then proves the general statement locally.
\end{example}

\begin{example}[Monomorphism stability]\label{ex:vi24-a11}
If $X\to Y$ is a monomorphism, then $X_{S'}\to Y_{S'}$ is a monomorphism because equality of maps can be tested after composing with the Cartesian projections.
\end{example}

\begin{example}[The diagonal commutes with base change]\label{ex:vi24-a12}
Using the universal property twice identifies $X_{S'}\times_{S'}X_{S'}$ with $(X\times_SX)_{S'}$ and the two diagonal maps correspond.
\end{example}

\begin{example}[Graphs commute with base change]\label{ex:vi24-a13}
Since graphs are pullbacks of diagonals, their compatibility with base change follows from associativity of pullbacks.
\end{example}

\begin{example}[A descent warning]\label{ex:vi24-a14}
A property visible after one base change need not hold before base change without a descent theorem. Geometric connectedness and geometric irreducibility are deliberately stronger notions than connectedness and irreducibility over the original field.
\end{example}
